\documentclass[10pt,a4paper]{article}

% -------------------------------------------------
% Packages
% -------------------------------------------------
\usepackage[a4paper,top=1.65cm,bottom=1.9cm,left=1.55cm,right=1.55cm,columnsep=0.65cm]{geometry}
\usepackage[T1]{fontenc}
\usepackage[utf8]{inputenc}
\usepackage{lmodern}
\usepackage{microtype}
\usepackage[table]{xcolor}
\usepackage{booktabs}
\usepackage{tabularx}
\usepackage{array}
\usepackage{graphicx}
\usepackage{caption}
\usepackage{enumitem}
\usepackage{fancyhdr}
\usepackage{titlesec}
\usepackage{hyperref}
\usepackage{amssymb}
\usepackage{pifont}
\usepackage[most]{tcolorbox}
\usepackage{multicol}
\usepackage{capt-of}

% -------------------------------------------------
% Colours and style inspired by the uploaded reference
% -------------------------------------------------
\definecolor{KentBlue}{HTML}{123B63}
\definecolor{DeepBlue}{HTML}{0B2545}
\definecolor{SoftBlue}{HTML}{EAF2FA}
\definecolor{TableHead}{HTML}{D9E7F5}
\definecolor{GoodGreen}{HTML}{1B7F5A}
\definecolor{WarnAmber}{HTML}{B7791F}
\definecolor{LightGrey}{HTML}{F6F8FA}
\definecolor{DarkGrey}{HTML}{343A40}

\hypersetup{
    colorlinks=true,
    linkcolor=KentBlue,
    urlcolor=KentBlue,
    citecolor=KentBlue,
    pdfauthor={Mukand Krishna},
    pdftitle={Individual Report: Classical Machine Learning Evaluation}
}

\pagestyle{fancy}
\fancyhf{}
\cfoot{\footnotesize\thepage}
\renewcommand{\headrulewidth}{0pt}
\setlength{\headheight}{13pt}
\raggedbottom

\titleformat{\section}{\large\bfseries\color{KentBlue}}{\thesection}{0.55em}{}
\titleformat{\subsection}{\normalsize\bfseries\color{DeepBlue}}{\thesubsection}{0.45em}{}
\titlespacing*{\section}{0pt}{8pt plus 2pt minus 1pt}{4pt}
\titlespacing*{\subsection}{0pt}{6pt plus 2pt minus 1pt}{3pt}

\captionsetup{hypcap=false,font=footnotesize,labelfont=bf,justification=centering}
\setlength{\parindent}{0pt}
\setlength{\parskip}{4pt}
\setlist[itemize]{leftmargin=1.2em,itemsep=1pt,topsep=2pt,parsep=0pt}

\newcolumntype{Y}{>{\raggedright\arraybackslash}X}
\newcolumntype{C}{>{\centering\arraybackslash}p{0.92cm}}
\newcommand{\tick}{\textcolor{GoodGreen}{\ding{51}}}
\newcommand{\warn}{\textcolor{WarnAmber}{\ding{115}}}
\newcommand{\separator}{%
    {\color{KentBlue}\rule{\textwidth}{1.1pt}}
}

\newtcolorbox{keybox}[1]{
    colback=SoftBlue,
    colframe=KentBlue,
    title=\textbf{#1},
    arc=1.5mm,
    boxrule=0.55pt,
    left=1.7mm,right=1.7mm,top=1mm,bottom=1mm,
    fonttitle=\footnotesize,
    fontupper=\footnotesize
}

\newtcolorbox{interpretbox}[1]{
    colback=LightGrey,
    colframe=DarkGrey,
    title=\textbf{#1},
    arc=1.5mm,
    boxrule=0.45pt,
    left=1.7mm,right=1.7mm,top=1mm,bottom=1mm,
    fonttitle=\footnotesize,
    fontupper=\footnotesize
}

\begin{document}

% Integrated first-page heading: no separate title page.
\begin{center}
    \vspace*{-0.35cm}
    {\LARGE\bfseries\color{KentBlue} AI and Cyber Security Project Individual Report}\\[2mm]
    {\normalsize \textbf{Mukand Krishna} \quad | \quad Login: \texttt{mk116}}\\[2mm]
    \separator
\end{center}


\vspace{2mm}
\begin{multicols}{2}

\section{Role and Technical Scope}

For my part, I worked on the classical machine learning experiments. I developed the grouped evaluation setup, implemented and tested heuristic baselines, trained Random Forest and RBF-kernel Support Vector Machin (SVM) models, and analysed whether the results depended too strongly on message-length-derived features.


I designed and applied a consistent evaluation framework across the heuristic baselines, Random Forest, and SVM experiments. This made the comparison fair because every method was assessed on the same pair-level dataset, the same grouped folds, and the same metric definitions.


The scope covered the full classical ML pipeline: preparing pair-level experimental data, applying grouped cross-validation, comparing non-ML baselines with trained models, and using no-length ablation to test shortcut dependence. This was important because strong model scores can be misleading if the task can be solved through simple artefacts in the data.

\section{Evaluation Design}

\subsection{Grouped 5-Fold Evaluation}

I implemented a grouped 5-fold evaluation strategy by dataset folder. In each fold, four complete dataset folders were used for training and one complete folder was held out for testing. This gave 88 training rows and 22 test rows per fold, with 8 positive training pairs and 2 positive test pairs.

This design was important because rows from the same dataset folder should not appear in both training and testing. Holding out complete folders reduced leakage risk and made the comparison between baselines, Random Forest, and SVM fairer.

\begin{center}
\captionof{table}{Grouped 5-fold evaluation structure.}
\label{tab:grouped-evaluation}
\renewcommand{\arraystretch}{1.15}
\footnotesize
\begin{tabularx}{\columnwidth}{lY}
\toprule
\rowcolor{TableHead}\textbf{Element} & \textbf{Description} \\
\midrule
Grouping & Dataset folder \\
Folds & 5 grouped folds \\
Training & 4 folders, 88 rows, 8 positives \\
Testing & 1 held-out folder, 22 rows, 2 positives \\
Purpose & Reduce leakage and keep comparison fair \\
\bottomrule
\end{tabularx}
\end{center}

\subsection{Evaluation Metrics}

The evaluation used both ranking and classification metrics. Top-1 Accuracy measured whether the correct pair was ranked first. Mean Reciprocal Rank (MRR) measured how highly the correct pair appeared in the ranked candidate list. F1 measured positive-pair classification quality, while ROC-AUC and PR-AUC assessed the trained models' discrimination performance.

\begin{center}
\captionof{table}{Metric interpretation.}
\label{tab:metrics}
\renewcommand{\arraystretch}{1.15}
\footnotesize
\begin{tabularx}{\columnwidth}{lY}
\toprule
\rowcolor{TableHead}\textbf{Metric} & \textbf{Meaning} \\
\midrule
Top-1 & True pair ranked first \\
MRR & True pair's rank quality \\
F1 & Balance between precision and recall \\
ROC-AUC & Separation across thresholds \\
PR-AUC & Positive-class ranking quality \\
\bottomrule
\end{tabularx}
\end{center}

\section{Heuristic Baselines}

Before training machine learning models, I tested simple heuristic baselines. This gave the project a reference point and helped show whether the trained models were learning something beyond simple rules. MRR was used to measure how highly the true pair was ranked within each scenario.

\end{multicols}
\begin{center}
\captionof{table}{Heuristic baseline experiment results.}
\label{tab:baseline-results}
\renewcommand{\arraystretch}{1.15}
\small
\begin{tabularx}{0.96\textwidth}{Y C C C Y}
\toprule
\rowcolor{TableHead}\textbf{Baseline} & \textbf{Top-1} & \textbf{MRR} & \textbf{F1} & \textbf{Interpretation} \\
\midrule
Highest message count & 0.00 & 0.31 & 0.00 & Volume alone failed \\
Longest duration & 0.30 & 0.58 & 0.30 & Duration alone was weak \\
Lowest mean length & 1.00 & 1.00 & 1.00 & Very strong, but a possible artefact \\
Highest evening ratio & 0.50 & 0.63 & 0.50 & Some signal, but not reliable \\
Highest short-message ratio & 1.00 & 1.00 & 1.00 & Strong length-related signal \\
Weighted contextual heuristic & 1.00 & 1.00 & 1.00 & Strong, but partly length-driven \\
\bottomrule
\end{tabularx}
\end{center}

\begin{center}
\includegraphics[width=0.84\textwidth]{figures/heuristic-baseline-top1.png}
\captionof{figure}{Top-1 accuracy of heuristic baselines using the original generated plot.}
\label{fig:baseline-top1}
\end{center}
\begin{multicols}{2}


The baseline results showed that message volume did not identify the true pair, while message-length-based rules performed extremely well. This helped motivate the no-length ablation later in the project.

Figure~\ref{fig:baseline-top1} shows why the baseline stage was necessary. Simple volume-based rules failed, whereas message-length-based rules performed perfectly. This provided clear evidence that the trained models had to be interpreted carefully, because excellent scores could partly reflect an exploitable synthetic pattern rather than generalisable behaviour.

\section{Random Forest and Support Vector Machine}

For the classical machine learning models, I trained Random Forest and RBF-kernel SVM models using the same grouped folds and metrics. Random Forest was useful because it can model non-linear feature interactions in a small tabular dataset. SVM was useful as a strong margin-based classifier, especially after feature scaling.

\begin{center}
\captionof{table}{Classical model results.}
\label{tab:model-results}
\begingroup
\renewcommand{\arraystretch}{1.25}
\scriptsize
\begin{tabularx}{\columnwidth}{Y C C C C C}
\toprule
\rowcolor{TableHead}\textbf{Model} & \textbf{T1} & \textbf{MRR} & \textbf{F1} & \textbf{ROC} & \textbf{PR} \\
\midrule
RF base & 1.00 & 1.00 & 1.00 & 1.00 & 1.00 \\
\addlinespace[2pt]
RF contextual & 1.00 & 1.00 & 1.00 & 1.00 & 1.00 \\
\addlinespace[2pt]
SVM base & 1.00 & 1.00 & 1.00 & 1.00 & 1.00 \\
\addlinespace[2pt]
SVM contextual & 1.00 & 1.00 & 1.00 & 1.00 & 1.00 \\
\bottomrule
\end{tabularx}
\endgroup
\end{center}

Both classical models achieved perfect performance with the base and contextual feature sets. This showed that the engineered pair-level features were highly predictive on PASYDA. However, the baseline results meant the scores still needed careful interpretation, because simple message-length rules could also perform extremely well.

\section{No-Length Ablation}

I also ran no-length ablation experiments by removing message-length-derived features and retraining the models with the same grouped folds. This tested whether the models were relying on a shortcut in the synthetic data.

\end{multicols}

\begin{center}
\captionof{table}{Full-feature versus no-length ablation results.}
\label{tab:ablation-results}
\begingroup
\renewcommand{\arraystretch}{1.22}
\small
\setlength{\tabcolsep}{4pt}
\begin{tabularx}{0.88\textwidth}{
  >{\raggedright\arraybackslash}p{0.17\textwidth}
  >{\centering\arraybackslash}p{0.095\textwidth}
  >{\centering\arraybackslash}p{0.115\textwidth}
  >{\centering\arraybackslash}p{0.090\textwidth}
  >{\centering\arraybackslash}p{0.095\textwidth}
  >{\centering\arraybackslash}p{0.115\textwidth}
  >{\centering\arraybackslash}p{0.090\textwidth}
}
\toprule
\rowcolor{TableHead}
\textbf{Experiment} & \textbf{Full T1} & \textbf{No-Length T1} & \textbf{T1 Drop} & \textbf{Full MRR} & \textbf{No-Length MRR} & \textbf{F1 Drop} \\
\midrule
Random Forest & 1.00 & 0.70 & 0.30 & 1.00 & 0.80 & 0.30 \\
\addlinespace[2pt]
SVM & 1.00 & 1.00 & 0.00 & 1.00 & 1.00 & 0.00 \\
\bottomrule
\end{tabularx}
\endgroup
\end{center}

\begin{center}
\includegraphics[width=0.68\textwidth]{figures/no-length-ablation-top1.png}
\captionof{figure}{Effect of removing message-length features using the original generated plot.}
\label{fig:no-length-ablation}
\end{center}
\begin{multicols}{2}


Figure~\ref{fig:no-length-ablation} highlights that Random Forest dropped from 1.00 to 0.70 Top-1 Accuracy when message-length features were removed, while SVM remained at 1.00. This suggests that Random Forest relied more strongly on length-based signals, whereas SVM was more robust under the same grouped evaluation setup.


\section{Deliverables Produced}

The following deliverables were produced as part of the classical machine learning work. Together, they provide experimental summaries, fold-level outputs, and written analysis needed to support reproducibility and interpretation.

\end{multicols}
\begin{center}
\captionof{table}{Classical machine learning deliverables.}
\label{tab:deliverables}
\renewcommand{\arraystretch}{1.18}
\small
\begin{tabularx}{0.96\textwidth}{>{\ttfamily\small\raggedright\arraybackslash}p{0.40\textwidth}Y}
\toprule
\rowcolor{TableHead}\normalfont\textbf{Deliverable} & \textbf{Purpose} \\
\midrule
baseline\_metrics\_summary.csv & Summary of heuristic baseline performance \\
model\_metrics\_summary.csv & Summary of RF, SVM, and MLP model performance \\
full\_vs\_no\_length\_comparison.csv & Direct comparison of full-feature and no-length experiments \\
baseline\_predictions\_by\_fold.csv & Fold-level baseline rankings for each candidate pair \\
model\_predictions\_by\_fold.csv & Fold-level model predictions and rankings \\
no\_length\_ablation\_experiment.md & Written analysis of shortcut dependence \\
\bottomrule
\end{tabularx}
\end{center}
\begin{multicols}{2}

\section{Critical Analysis}

The most important lesson I got from this work is that strong model scores are not enough by themselves. At first, the perfect results looked like a simple success. However, the baselines showed that a hand-made message-length rule could also solve the task. That changed how I interpreted the models.

The baseline experiments and no-length ablation showed that perfect performance could be partly driven by message-length artefacts in the synthetic dataset. This made the final interpretation more cautious, transparent, and defensible.

\section{Limitations and Future Work}

Although the grouped evaluation reduced leakage risk, the experiments were still conducted on a synthetic dataset. Therefore, the results should be validated on larger and more diverse data before making strong claims about real-world generalisation.

Future work should include stronger calibration analysis, more detailed threshold tuning, external validation on a larger dataset, deeper fold-level error analysis, and comparison with additional models under the same grouped evaluation protocol.

\section{Conclusion}


The classical machine learning experiments demonstrated strong predictive performance, but the most valuable contribution was the evaluation framework and the cautious interpretation. By combining grouped folds, heuristic baselines, and no-length ablation, the work showed both what the models achieved and why their perfect scores should be interpreted carefully.

\end{multicols}
\end{document}
